% GENERATED FILE. Edit canonical lessons, then run npm run generate:books.
% canonical-source-hash: fnv1a64:272b75d22310dce6
% canonical-lessons: ZH-C370-yinsu, ZH-C370-bijing, ZH-C370-guanjian, ZH-C370-zuoyong, ZH-C370-zongzhi

\chapter{The Key Point}
\label{ch:zh-the-key-point}
\hlchaptermodality{370}

\begin{chapteropening}
\textbf{By the end of this chapter:} \emph{I can name a factor and the key point, say what effect something had, and sum up.}

Complete the last lesson of chapter 370: I can name a factor and the key point, say what effect something had, and sum up.
\end{chapteropening}

\section[yīnsù]{yīnsù — factor}

\label{lesson:ZH-C370-yinsu}



\begin{quote}
Before the new one: say the Mandarin for secondly, then the Mandarin for in addition.
\end{quote}

\subsection*{You'll want to know: yīnsù}
\emph{yīnsù} — \textquotedblleft{}factor\textquotedblright{}.

\begin{quote}
\emph{wǒ juéde tiānqì shì yí ge zhòngyào de yīnsù} — I think the weather is an important factor.
\end{quote}

Its characters come later; for now it is a word for the ear.

\begin{cousinweb}
\emph{yīnsù} is built from \zh{因} 'a cause' and \emph{sù}, both inherited from Middle Chinese, \emph{'jin} and \emph{suH} (Guangyun, in Baxter's transcription), which glosses the first 'to rely on' and 'a cause', and the second 'empty' and 'silk'. How the pair came to mean a factor is not settled.
\end{cousinweb}

\begin{grammarlens}[title={What you've built}]
One more word for giving reasons.
\end{grammarlens}

\subsection*{Your turn}
\begin{itemize}
\raggedright
  \item \emph{Say it:} \emph{yīnsù}
  \item \emph{Say it:} \emph{yīnsù}, once more
  \item \emph{Say it:} \emph{wǒ juéde tiānqì shì yí ge zhòngyào de yīnsù} — I think the weather is an important factor.
  \item \emph{Recall:} say the Mandarin for secondly, then the Mandarin for in addition, then say \emph{yīnsù} again
\end{itemize}

\begin{tcolorbox}[breakable,colback=teal!4,colframe=teal!35!black,title={Before you move on}]
What does yīnsù mean? (\textquotedblleft{}Factor\textquotedblright{}.) Say it once more.
\end{tcolorbox}
\section[bìjìng]{bìjìng — after all; all things considered}

\label{lesson:ZH-C370-bijing}



\begin{quote}
Before the new one: say the Mandarin for in addition, then the Mandarin for a factor.
\end{quote}

\subsection*{You'll want to know: bìjìng}
\emph{bìjìng} — \textquotedblleft{}after all; all things considered\textquotedblright{}.

\begin{quote}
\emph{tā bìjìng shì háizi} — He's only a child, after all.
\end{quote}

Its characters come later; for now it is a word for the ear.

\begin{cousinweb}
\emph{bìjìng} is built from \emph{bì} and \emph{jìng}, both inherited from Middle Chinese, \emph{pjit} and \emph{kjaengH} (Guangyun, in Baxter's transcription), which glosses the first 'the end' and the second 'to use up' and 'the end'.
\end{cousinweb}

\begin{grammarlens}[title={What you've built}]
One more word for giving reasons.
\end{grammarlens}

\subsection*{Your turn}
\begin{itemize}
\raggedright
  \item \emph{Say it:} \emph{bìjìng}
  \item \emph{Say it:} \emph{bìjìng}, once more
  \item \emph{Say it:} \emph{tā bìjìng shì háizi} — He's only a child, after all.
  \item \emph{Recall:} say the Mandarin for in addition, then the Mandarin for a factor, then say \emph{bìjìng} again
\end{itemize}

\begin{tcolorbox}[breakable,colback=teal!4,colframe=teal!35!black,title={Before you move on}]
What does bìjìng mean? (\textquotedblleft{}After all; all things considered\textquotedblright{}.) Say it once more.
\end{tcolorbox}
\section[guānjiàn]{guānjiàn — key; crucial point}

\label{lesson:ZH-C370-guanjian}



\begin{quote}
Before the new one: say the Mandarin for a factor, then the Mandarin for all things considered.
\end{quote}

\subsection*{You'll want to know: guānjiàn}
\emph{guānjiàn} — \textquotedblleft{}key; crucial point\textquotedblright{}.

\begin{quote}
\emph{wǒ juéde guānjiàn shì shíjiān} — I think the key thing is time.
\end{quote}

Its characters come later; for now it is a word for the ear.

\begin{cousinweb}
\emph{guānjiàn} is built from \zh{关} 'to shut' and \emph{jiàn}, both inherited from Middle Chinese, \emph{kwaen} and \emph{gjen} (Guangyun, in Baxter's transcription), which glosses the first, quoting an older dictionary, 'a bar across a door', and the second 'a lock'. Middle Chinese \emph{gjen} would regularly give qián; the dictionary also lists the character in the rising tone, which fits jiàn.
\end{cousinweb}

\begin{grammarlens}[title={What you've built}]
One more word for giving reasons.
\end{grammarlens}

\subsection*{Your turn}
\begin{itemize}
\raggedright
  \item \emph{Say it:} \emph{guānjiàn}
  \item \emph{Say it:} \emph{guānjiàn}, once more
  \item \emph{Say it:} \emph{wǒ juéde guānjiàn shì shíjiān} — I think the key thing is time.
  \item \emph{Recall:} say the Mandarin for a factor, then the Mandarin for all things considered, then say \emph{guānjiàn} again
\end{itemize}

\begin{tcolorbox}[breakable,colback=teal!4,colframe=teal!35!black,title={Before you move on}]
What does guānjiàn mean? (\textquotedblleft{}Key; crucial point\textquotedblright{}.) Say it once more.
\end{tcolorbox}
\section[\texorpdfstring{\zh{作用} (zuòyòng)}{作用 (zuòyòng)}]{\zh{作用} (zuòyòng) — effect; function}

\label{lesson:ZH-C370-zuoyong}



\begin{quote}
Before the new one: say the Mandarin for all things considered, then the Mandarin for the crucial point.
\end{quote}

\subsection*{You'll want to know: \zh{作用}}
\textbf{\zh{作用}} — \emph{zuòyòng} — \textquotedblleft{}effect; function\textquotedblright{}.

\begin{quote}
\textbf{\zh{我说的话没有作用}} \emph{wǒ shuō de huà méiyǒu zuòyòng} — What I said had no effect.
\end{quote}

\begin{cousinweb}
\textbf{\zh{作用}} (\emph{zuòyòng}) is built from \zh{作} 'to do' and \zh{用} 'to use', both inherited from Middle Chinese, \emph{tsak} and \emph{yowngH} (Guangyun, in Baxter's transcription), which glosses the first 'to do' and 'to rise', and the second 'to use'.
\end{cousinweb}

\begin{grammarlens}[title={What you've built}]
One more word for giving reasons.
\end{grammarlens}

\subsection*{Your turn}
\begin{itemize}
\raggedright
  \item \emph{Say it:} \emph{zuòyòng}
  \item \emph{Say it:} \emph{zuòyòng}, once more
  \item \emph{Say it:} \emph{wǒ shuō de huà méiyǒu zuòyòng} — What I said had no effect.
  \item \emph{Recall:} say the Mandarin for all things considered, then the Mandarin for the crucial point, then say \emph{zuòyòng} again
\end{itemize}

\begin{tcolorbox}[breakable,colback=teal!4,colframe=teal!35!black,title={Before you move on}]
What does \zh{作用} mean? (\textquotedblleft{}Effect; function\textquotedblright{}.) Say it once more.
\end{tcolorbox}
\section[zǒngzhī]{zǒngzhī — in short; in a word}

\label{lesson:ZH-C370-zongzhi}



\begin{quote}
Before the new one: say the Mandarin for the crucial point, then the Mandarin for the effect something has.
\end{quote}

\subsection*{You'll want to know: zǒngzhī}
\emph{zǒngzhī} — \textquotedblleft{}in short; in a word\textquotedblright{}.

\begin{quote}
\emph{zǒngzhī wǒ juéde nǐ yīnggāi qù} — In short, I think you should go.
\end{quote}

Its characters come later; for now it is a word for the ear.

\begin{cousinweb}
\emph{zǒngzhī} is built from \emph{zǒng} and \emph{zhī}, both inherited from Middle Chinese, \emph{tsuwngX} and \emph{tsyi} (Guangyun, in Baxter's transcription), which glosses the first 'to bind together' and 'all', and the second 'to go'. How the pair came to mean in short is not settled.
\end{cousinweb}

\begin{grammarlens}[title={What you've built}]
One more word for giving reasons.
\end{grammarlens}

\subsection*{Your turn}
\begin{itemize}
\raggedright
  \item \emph{Say it:} \emph{zǒngzhī}
  \item \emph{Say it:} \emph{zǒngzhī}, once more
  \item \emph{Say it:} \emph{zǒngzhī wǒ juéde nǐ yīnggāi qù} — In short, I think you should go.
  \item \emph{Recall:} say the Mandarin for the crucial point, then the Mandarin for the effect something has, then say \emph{zǒngzhī} again
\end{itemize}

\begin{tcolorbox}[breakable,colback=teal!4,colframe=teal!35!black,title={Before you move on}]
What does zǒngzhī mean? (\textquotedblleft{}In short; in a word\textquotedblright{}.) Say it once more.
\end{tcolorbox}
